% =============================================================================
% IEEEtran Template for DS108 Electrimight Report
% =============================================================================
\documentclass[conference]{IEEEtran}
\IEEEoverridecommandlockouts

% ---------- Packages ----------
\usepackage{cite}
\usepackage{amsmath,amssymb,amsfonts}
\usepackage{algorithmic}
\usepackage{graphicx}
\usepackage{textcomp}
\usepackage{xcolor}
\usepackage{booktabs}
\usepackage{multirow}
\usepackage{pgfplots}
\pgfplotsset{compat=1.17}
\usepackage{subcaption}

% ---------- Hyphenation ----------
\hyphenation{op-tical net-works semi-conduc-tor}

% ---------- Document ----------
\begin{document}

% ---------- Title ----------
\title{Wavelet-Based Feature Engineering and GAN Augmentation 
for Steel Industry Energy Consumption Dataset Construction}

% ---------- Authors ----------
\author{
    \IEEEauthorblockN{Author Name\textsuperscript{1}, 
    Author Name\textsuperscript{2}, Author Name\textsuperscript{3}}
    \IEEEauthorblockA{\textsuperscript{1}\textit{Department of Data Science}, 
    University Name, City, Country}
    \IEEEauthorblockA{\textsuperscript{2}\textit{Department of Electrical Engineering}, 
    University Name, City, Country}
    \IEEEauthorblockA{\textsuperscript{3}\textit{Department of Computer Science}, 
    University Name, City, Country}
    \IEEEauthorblockA{Email: author1@university.edu}
}

\maketitle

% ---------- Abstract ----------
\begin{abstract}
Industrial electricity consumption forecasting and risk 
classification require high-quality datasets with rich features. 
This paper presents a comprehensive preprocessing pipeline for 
the Steel Industry Energy Consumption dataset, a real-world 
time-series dataset collected from a South Korean steel plant. 
We propose a multi-domain feature engineering framework that 
extracts time-domain statistics (lag and rolling features), 
frequency-domain characteristics via Discrete Wavelet Transform 
(DWT) with Daubechies-4 wavelets, and physical-domain derivatives 
including Apparent Power and Phase Angle. Furthermore, anomaly 
labels for idling, energy leakage, and local overload are 
generated using rule-based physical thresholds. To address class 
imbalance, a Generative Adversarial Network (GAN) is employed 
for synthetic data augmentation. The processed dataset expanded 
from 7 original features to over 40 engineered features with 
three binary anomaly indicators. Experimental validation confirms 
that the extracted wavelet and physical features effectively 
capture non-stationary load behaviors and latent fault 
signatures, providing a robust foundation for downstream 
forecasting and anomaly detection models.
\end{abstract}

\begin{IEEEkeywords}
Energy consumption, data preprocessing, discrete wavelet 
transform, feature engineering, generative adversarial network, 
anomaly detection, steel industry.
\end{IEEEkeywords}

% =============================================================================
\section{INTRODUCTION}
\label{sec:introduction}
% =============================================================================

The steel industry accounts for approximately 7--9\% of global 
anthropogenic CO$_2$ emissions, with electricity consumption 
representing a dominant operational cost. Accurate load 
forecasting and early detection of electrical anomalies are 
therefore critical for energy optimization and predictive 
maintenance in steel manufacturing plants.

Despite the availability of public industrial energy datasets, 
existing corpora typically provide only raw meter readings 
without derived physical indicators or anomaly annotations. 
This limitation constrains the applicability of supervised 
learning models for fault diagnosis. Furthermore, anomalous 
events---such as idling motors, insulation degradation, and 
localized overloads---are rare, resulting in severe class 
imbalance that degrades classifier performance.

To address these limitations, this paper proposes a systematic 
preprocessing and dataset construction pipeline for steel 
industry electricity data. The contributions are threefold: 
(i) a multi-domain feature engineering framework that augments 
raw signals with time-lag statistics, Discrete Wavelet Transform 
(DWT) coefficients, and physically meaningful derivatives 
including Apparent Power and Phase Angle; (ii) a 
physics-informed anomaly labeling scheme that identifies 
idling, energy leakage, and local overload events; and (iii) 
a Generative Adversarial Network (GAN) for synthetic 
minority-class augmentation to mitigate data imbalance.

The remainder of this paper is organized as follows. 
Section~\ref{sec:related} reviews related work. 
Section~\ref{sec:methodology} details the proposed methodology. 
Section~\ref{sec:results} presents results and discussion. 
Finally, Section~\ref{sec:conclusion} concludes the paper.

% =============================================================================
\section{RELATED WORK}
\label{sec:related}
% =============================================================================

\subsection{Industrial Energy Consumption Forecasting}

Industrial load forecasting has been extensively studied using 
statistical and machine learning approaches. Auto-regressive 
integrated moving average (ARIMA) models provide interpretable 
baselines but fail to capture nonlinear dynamics~\cite{hippert2001}. 
More recently, deep learning architectures such as Long Short-Term 
Memory (LSTM) networks have demonstrated superior accuracy in 
modeling temporal dependencies~\cite{bouktif2020}. However, these 
studies predominantly rely on aggregated system-level data, whereas 
end-user industrial loads---particularly in steel manufacturing---
exhibit abrupt fluctuations driven by furnace scheduling and 
rolling mill operations~\cite{khuntia2016}.

\subsection{Wavelet Transform in Time-Series Analysis}

Frequency-domain analysis complements time-domain methods by 
revealing transient behaviors invisible to rolling statistics. 
The Fast Fourier Transform (FFT) decomposes signals into 
frequency components but sacrifices temporal localization. 
In contrast, the Discrete Wavelet Transform (DWT) employs 
scaled and translated mother wavelets to capture both frequency 
and time information, making it particularly suitable for 
non-stationary industrial load signals~\cite{mallat1989}. 
Zhang et al.~\cite{zhang2024} demonstrated that a DWT-LSTM 
hybrid model achieves MAPE as low as 0.59\% for short-term 
load forecasting.

\subsection{GAN and Synthetic Data Augmentation}

Class imbalance poses a persistent challenge in anomaly 
detection, where abnormal events constitute less than 1\% of 
observations~\cite{chandola2009}. Generative Adversarial Networks 
(GANs), introduced by Goodfellow et al.~\cite{goodfellow2014}, 
have emerged as a powerful paradigm for synthetic data generation. 
For time-series applications, Yoon et al.~\cite{yoon2019} proposed 
TimeGAN, which combines adversarial training with supervised 
sequence modeling to preserve temporal dynamics.

\begin{table}[htbp]
\caption{COMPARISON OF RELATED WORKS}
\label{tab:related}
\centering
\begin{tabular}{@{}ccccccc@{}}
\toprule
Ref & Forecast & DWT & Physical & Anomaly & GAN & Domain \\
\midrule
\cite{hippert2001} & Yes & No & No & No & No & General \\
\cite{zhang2024} & Yes & Yes & No & No & No & Power Grid \\
\cite{tang2025} & No & No & No & No & Yes & Manufacturing \\
Ours & Yes & Yes & Yes & Yes & Yes & Steel \\
\bottomrule
\end{tabular}
\end{table}

% =============================================================================
\section{METHODOLOGY}
\label{sec:methodology}
% =============================================================================

\subsection{Dataset Overview}

The dataset used in this study is the Steel Industry Energy 
Consumption dataset~\cite{uci2023}, publicly available from 
the UCI Machine Learning Repository. It contains 35,040 
observations recorded at 15-minute intervals from a South 
Korean steel manufacturing plant.

\subsection{Data Preprocessing}

The preprocessing stage comprises data inspection, duplicate 
removal, and missing value imputation. For time-series data, 
linear interpolation is applied:
\begin{equation}
y(t) = y(t_1) + \frac{y(t_2) - y(t_1)}{t_2 - t_1}(t - t_1)
\label{eq:interp}
\end{equation}

\subsection{Time-Domain Feature Engineering}

Lag features are constructed with delays $k \in \{1, 2, 4, 96\}$ 
corresponding to 15 min, 30 min, 1 h, and 24 h:
\begin{equation}
x_{\text{lag}}(k) = x(t - k)
\label{eq:lag}
\end{equation}

Rolling window statistics include mean, standard deviation, 
and skewness over windows $w \in \{24, 48, 96\}$.

Cyclical encoding of the Number of Seconds from Midnight (NSM) 
ensures continuity at midnight:
\begin{align}
\text{NSM}_{\sin} &= \sin\left(\frac{2\pi \cdot \text{NSM}}{86400}\right) \\
\text{NSM}_{\cos} &= \cos\left(\frac{2\pi \cdot \text{NSM}}{86400}\right)
\end{align}

\subsection{Frequency-Domain Feature Extraction (DWT)}

The Discrete Wavelet Transform decomposes the signal into 
approximation $cA$ and detail $cD$ coefficients:
\begin{equation}
S(t) = cA_j(t) + \sum_{i=1}^{j} cD_i(t)
\label{eq:dwt}
\end{equation}
We employ the Daubechies-4 (db4) wavelet with 3 decomposition 
levels. A rolling window of 64 samples (16 hours) is used to 
compute wavelet features at each time step.

\subsection{Physical-Domain Feature Engineering}

Apparent Power $S$ and Phase Angle $\varphi$ are computed as:
\begin{align}
S &= \sqrt{P^2 + Q^2} \\
\varphi &= \arccos(\text{PF})
\end{align}
where $P$ is active power (Usage\_kWh), $Q$ is reactive power, 
and PF is the power factor clipped to $[0, 1]$.

\subsection{Anomaly Labeling Strategy}

Three anomaly types are defined using physics-informed thresholds:
\begin{itemize}
\item \textbf{Idling:} Light load during off-hours with high 
consumption and PF $< 0.5$
\item \textbf{Leakage:} Rolling mean exceeds baseline by 5\%
\item \textbf{Overload:} Usage exceeds $Q_3 + 3\times$IQR with 
high reactive power and PF $< 0.7$
\end{itemize}

\subsection{GAN-Based Data Augmentation}

A fully connected GAN is trained to generate synthetic 
minority-class samples. The generator and discriminator 
architectures follow standard deep fully connected networks 
with LeakyReLU activations. The minimax objective is:
\begin{equation}
\min_G \max_D \mathcal{L}(D,G) = \mathbb{E}_{x}[\ln D(x)] + 
\mathbb{E}_{z}[\ln(1-D(G(z)))]
\label{eq:gan}
\end{equation}

% =============================================================================
\section{RESULTS AND DISCUSSION}
\label{sec:results}
% =============================================================================

\subsection{Data Profiling}

[Insert descriptive statistics table and time-series plots]

\subsection{Feature Engineering Evaluation}

[Insert correlation heatmap, DWT decomposition figure, 
and S-$\varphi$ scatter plot]

\subsection{Anomaly Distribution}

[Insert anomaly label distribution table and timeline plot]

\subsection{GAN Augmentation Quality}

[Insert real vs. synthetic comparison table and PCA/t-SNE plots]

% =============================================================================
\section{CONCLUSION}
\label{sec:conclusion}
% =============================================================================

This paper presented a comprehensive preprocessing and dataset 
construction pipeline for the Steel Industry Energy Consumption 
dataset. By integrating time-domain statistics, Discrete Wavelet 
Transform coefficients, and physics-informed derivatives, the 
proposed framework enriched the original dataset from 11 to 
over 40 features. Three anomaly classes were labeled using 
rule-based physical thresholds, and a GAN-based augmentation 
module synthesized 500 minority-class samples.

The key contributions are:
\begin{itemize}
\item A multi-domain feature engineering framework extracting 
temporal, spectral, and physical characteristics
\item A physics-informed anomaly labeling strategy detecting 
idling, leakage, and overload without manual annotation
\item A GAN augmentation pipeline with first-order statistical 
fidelity validated via PCA and t-SNE
\end{itemize}

Future work will explore adaptive thresholding via clustering, 
TimeGAN architectures for better temporal preservation, and 
downstream evaluation using LSTM and Transformer models.

% =============================================================================
% REFERENCES
% =============================================================================
\begin{thebibliography}{00}

\bibitem{hippert2001}
H.~S. Hippert, C.~E. Pedreira, and R.~C. Souza, ``Neural networks 
for short-term load forecasting: A review and evaluation,'' 
\emph{IEEE Trans. Power Syst.}, vol.~16, no.~1, pp. 44--55, 
Feb. 2001.

\bibitem{bouktif2020}
S.~Bouktif \emph{et al.}, ``Optimal deep learning LSTM model for 
electric load forecasting using feature selection and genetic 
algorithm,'' \emph{IEEE Access}, vol.~8, pp. 152,431--152,445, 2020.

\bibitem{khuntia2016}
S.~R. Khuntia, J.~L. Rueda, and M.~A.~M.~M. van der Meijden, 
``Forecasting the load of electrical power systems in mid- and 
long-term horizons: A review,'' \emph{IET Gener. Transm. Distrib.}, 
vol.~10, no.~16, pp. 3971--3977, 2016.

\bibitem{uci2023}
``Steel Industry Energy Consumption Dataset,'' UCI Machine 
Learning Repository. [Online]. Available: 
https://archive.ics.uci.edu/dataset/851/

\bibitem{zhang2024}
N.~Zhang \emph{et al.}, ``Deep learning modeling in electricity 
load forecasting,'' \emph{Energy Rep.}, vol.~11, pp. 3544--3556, 
2024.

\bibitem{mallat1989}
S.~G. Mallat, ``A theory for multiresolution signal decomposition: 
The wavelet representation,'' \emph{IEEE Trans. Pattern Anal. 
Mach. Intell.}, vol.~11, no.~7, pp. 674--693, Jul. 1989.

\bibitem{chandola2009}
V.~Chandola, A.~Banerjee, and V.~Kumar, ``Anomaly detection: 
A survey,'' \emph{ACM Comput. Surv.}, vol.~41, no.~3, pp. 1--58, 
2009.

\bibitem{goodfellow2014}
I.~J. Goodfellow \emph{et al.}, ``Generative adversarial nets,'' 
in \emph{Advances in Neural Information Processing Systems 
(NeurIPS)}, Montreal, Canada, 2014, pp. 2672--2680.

\bibitem{yoon2019}
J.~Yoon, D.~Jarrett, and M.~van der Schaar, ``Time-series 
generative adversarial networks,'' in \emph{Proc. 33rd 
Conference on Neural Information Processing Systems (NeurIPS)}, 
Vancouver, Canada, 2019, pp. 5509--5519.

\bibitem{tang2025}
P.~Tang \emph{et al.}, ``Time series data augmentation for 
energy consumption data based on improved TimeGAN,'' 
\emph{Sensors}, vol.~25, no.~2, p.~493, Jan. 2025.

\end{thebibliography}

\end{document}
